% Generated draft: requires mathematical and source review.
\documentclass[11pt]{article}
\usepackage[margin=25mm]{geometry}
\usepackage{fontspec}
\setmainfont{Latin Modern Roman}
\usepackage{amsmath,amssymb,mathtools,unicode-math}
\setmathfont{Latin Modern Math}
\usepackage{enumitem,needspace,xurl,hyperref}
\setlength{\parindent}{0pt}
\setlength{\parskip}{5pt}
\setlength{\emergencystretch}{3em}
\newcommand{\N}{\mathbb{N}}
\newcommand{\Z}{\mathbb{Z}}
\newcommand{\Q}{\mathbb{Q}}
\newcommand{\R}{\mathbb{R}}
\newcommand{\C}{\mathbb{C}}
\newcommand{\F}{\mathbb{F}}
\newcommand{\A}{\mathbb{A}}
\newcommand{\PP}{\mathbb{P}}
\newcommand{\E}{\mathbb{E}}
\newcommand{\eps}{\varepsilon}
\newcommand{\dd}{\,\mathrm{d}}
\begin{document}
\section*{UnsolvedMath ID: 7000030}
\textbf{Source problem code:} catalogue entry
\subsection{Definitions and mathematical statement}
Let $M\subset\R^3$ be a complete noncompact convex surface in the source's intended sense, with regularity at least $C^2$ where its principal curvatures $k_1(p),k_2(p)$ are defined. Completeness refers to the induced Riemannian metric. The absolute difference of the curvatures is independent of the choice of unit normal and of their ordering.

The assertion is
\[
 \inf_{p\in M}|k_1(p)-k_2(p)|=0,
\]
or equivalently
\[
 \forall\varepsilon>0\quad\exists p\in M,
 \qquad |k_1(p)-k_2(p)|<\varepsilon.
\]
It rules out a uniform positive lower bound on the curvature difference. It does not claim existence of an actual umbilic point, attainment of the infimum, local spherical geometry, or that the approximating points must escape to infinity.

The word convex is crucial but is not defined in the record. Being the boundary of a convex body, having a semidefinite second fundamental form, strict convexity, positive Gaussian curvature, and mean convexity are different conditions. Cylindrical or other developable cases make the weak-versus-strict distinction material, so the target must be checked against the source before adopting one convention. The research summary's identification of convexity with mean/positive curvature and its description as an entire graph are not justified hypotheses of the extracted statement. No strict-curvature or graph assumption is silently added.
\subsection{Short English statement}
Must a complete noncompact convex surface have points with arbitrarily small difference between its principal curvatures? This is an approximate-umbilicity assertion, not a claim that an umbilic point exists. The intended strictness of convexity is essential.

The supplied summary reports no definitive general proof and labels the problem open, but its statements about actual umbilics, local spherical shape, mean convexity, and entire graphs overstate or change the supplied target. They are retained as formulation cautions. The OPEN-TRIAGE label and 2019 survey are supplied assessments, with current status marked for review.
\subsection{Sources}
\begin{itemize}
\item Local UnsolvedMath registry, permanent ID 7000030, source code catalogue entry, ``Geometry of Curves and Surfaces --- Problem 8.3,'' proposed by Mohammad Ghomi, 2019. Dataset curation and metadata: Ulam AI and UnsolvedMath Contributors, CC BY 4.0.
\item Mohammad Ghomi, \emph{Open Problems in Geometry of Curves and Surfaces} (2019), Problem 8.3, PDF page 16, \url{https://people.math.gatech.edu/~ghomi/Papers/op.pdf}; normalized pdf-text extraction, accessed 2026-07-29.
\item No later precise theorem or primary resolution reference is supplied. Source-specific provenance: the author survey labels the item a problem; current status and statement-reuse terms NEEDS\_REVIEW for the public author-hosted PDF, separately from the dataset licence.
\end{itemize}
\end{document}
